\subsection{Desarrollo}



El desarrollo del proyecto se ha estructurado en etapas que
responden a los objetivos planteados y al cronograma
establecido.


\paragraph{Análisis del modelo base.}

Se realizó ingeniería inversa del código original en
\textsc{Mathematica} proporcionado por \cite{Damian_2022},
identificando la función potencial, sus descriptores
(cuadrado del gradiente, matriz hessiana, autovalores) y la
lógica del método de optimización empleado originalmente
(templado simulado con refinamiento por gradiente conjugado).
Este análisis permitió extraer las expresiones simbólicas que
rigen el modelo.


\paragraph{Implementación simbólica en Python.}

Con base en las expresiones obtenidas, se desarrolló un
módulo en Python que utiliza \texttt{sympy} para definir la
función potencial en sus dos versiones (con y sin el término
de lifting asociado a las $\widehat{D5}$-branas). El módulo
calcula también el cuadrado del gradiente, la matriz hessiana y
sus autovalores, proporcionando así todos los elementos
necesarios para evaluar la calidad de las soluciones y
aplicar los criterios de penalización durante la optimización.


\paragraph{Validación de la implementación.}

Para garantizar la correcta interpretación del modelo, se
compararon manualmente las expresiones simbólicas generadas
por el código en Python con las del código original en
\textsc{Mathematica}. La correspondencia obtenida para la
función potencial y los elementos de la matriz jacobiana
confirma que la implementación es fiel al modelo de
referencia.


\paragraph{Experimentación preliminar con algoritmos
genéticos.}

Una vez validada la implementación base, se sustituyó el
método de optimización original por un algoritmo genético
simple (\texttt{SGA}). Utilizando la misma base de datos de
configuraciones iniciales que en \cite{Damian_2022}, se
obtuvieron resultados preliminares que incluyen gráficas de
convergencia y la clasificación de las soluciones en estables
o inestables (según la presencia de taquiones) y en dS o AdS
(según el signo de la energía potencial). Estos resultados
sirven como punto de partida para la comparación con otros
métodos. Además de haber arrojado elementos interesantes
y relevantes a la búsqueda.


\paragraph{Implementación de los métodos restantes.}

Con la infraestructura de software ya establecida, se
procederá a implementar secuencialmente los siguientes
algoritmos metaheurísticos, tal como se indica en el
cronograma: optimización por enjambre de partículas (PSO),
algoritmos de estimación de distribución (EDAs), optimización
por colonia de hormigas (ACO), optimización por maleza
invasora (IWO), búsqueda de patrones (VCS) y búsqueda por
sintonización adaptativa (ASO). Cada método se evaluará sobre
las mismas bases de datos para garantizar la comparabilidad
de los resultados.


\paragraph{Análisis comparativo.}

Finalmente, se realizará un análisis comparativo del
desempeño de cada metaheurística en términos de eficiencia
(número de evaluaciones necesarias para converger), calidad
de las soluciones (valor de la función objetivo en el mínimo
encontrado) y capacidad de descubrimiento (cantidad de vacíos
estables identificados). Los resultados permitirán determinar
cuáles métodos se adaptan mejor a la geometría del problema y
si es posible incrementar el número de configuraciones
estables reportadas respecto al trabajo de referencia.
